\section{A twelvefold with two active adjoint embeddings}\label{sec:two-active}

\begin{proposition}\label{prop:two-active-twelvefold}
There exists a polarized abelian twelvefold over a number field with
full geometric rational endomorphism algebra a sextic CM field and
relative signatures \((2,2),(2,2),(0,4)\), up to orientation.
Its adjoint Mumford--Tate group is rationally simple with six
geometric factors of type \(A_1\), exactly two of them Hodge-active.
It has density-one ordinary reduction after finite extension.
\end{proposition}

\begin{proof}
Let \(R=\Q(\zeta_{13})^+\) and let \(F\) be its cubic subfield.
Number the six real embeddings of \(R\) cyclically from zero to five,
so the fibers over the embeddings of \(F\) are \(\{j,j+3\}\) for
\(j=0,1,2\). Take the quaternion algebra \(D/R\) split at the real
embeddings numbered zero and one and at every finite place, and
ramified at the other four real embeddings. Its existence follows
from quaternion reciprocity \cite[Theorem~14.6.1]{voight}.
Put \(E=F(\sqrt{-1})\). The corestriction
\(\operatorname{Cor}_{R/F}D\) is a central simple algebra of degree
four. After extension to \(E\) it is split locally everywhere and
hence split globally \cite[\S14.6]{voight}. On a simple module
\(W=E^4\), the group \(\operatorname{Res}_{R/\Q}\mathrm{SL}_1(D)\)
acts by corestriction. Over an algebraic closure the block for an
embedding of \(E\) above \(j\) is the tensor of the two standard
representations indexed by \(\{j,j+3\}\). This is the corestriction
representation of \cite[Example~3.1]{moonenclassification}, with
the degree-two fibers taken over \(F\). Adjoin the commuting scalar
torus of elements of \(\operatorname{Res}_{E/\Q}\mathbb G_m\) whose
product with their CM conjugate lies in \(\mathbb G_m\). Let \(G_0\)
be the image of the resulting product on \(W\).

The real Hodge action of \(z\in\mathbb C^\times\) at the split
quaternion embeddings is rotation by \(z/|z|\) on the real standard
module. It is trivial on the compact quaternion factors. Multiply
this action by the scalar \(|z|\) at the two embedding pairs
containing a split factor and by the complex scalar \(z\) at the
remaining pair. The scalar factors have the common norm \(|z|^2\).
On each mixed real block the resulting action is multiplication by
the real standard matrix of \(z\) in the split tensor factor. On the
definite block it is complex scalar multiplication. These formulas
define a real algebraic Hodge homomorphism of weight one with the
required signatures.

On a mixed block take the real part of the skew-Hermitian form given
by the split alternating form tensored with an invariant positive
Hermitian form on the compact standard module. On the definite block
use the imaginary part of an invariant positive Hermitian form.
Choose the signs so the Hodge complex structures are positive.
The derived group preserves these forms, the scalar torus acts with
its common norm as similitude, and the adjoint of the \(E\)-action is
CM conjugation. These conditions on an alternating form are rational
linear conditions. Rational approximation within their real solution
space supplies a rational form satisfying the same open positivity
condition. A lattice scaled to make the form integral then gives a
polarized complex abelian variety.

Vary the Hodge homomorphism in its connected real conjugacy class.
A rational adjoint subgroup containing the entire class contains
both split simple factors. Indeed, conjugating in a split factor and
taking ratios of the conjugates generates that factor, by simplicity
and the nontrivial projection of the Hodge action. Rationality and
Galois transitivity then give all six factors. There are countably
many proper rational adjoint subgroups, and the homomorphisms in the
class lying in any one of them form a proper real analytic locus.
Thus a very general member has full adjoint group \(G_0^{\mathrm{ad}}\).
The commutant of its Mumford--Tate representation is exactly \(E\).
Indeed the three distinct tensor supports distinguish blocks over
different embeddings of \(F\). Between the two orientations of the
definite pair the scalar Hodge actions differ, so there is no
intertwiner. Rationality and Galois transitivity on the pairs give
the same conclusion for the two mixed pairs. The commutant therefore
acts by independent scalars on the six blocks over an algebraic
closure, and its rational form is the given copy of \(E\).

Noot's specialization theorem gives a number-field point in the same
connected Shimura component with this Mumford--Tate representation
\cite[Theorem~1.7 and its proof]{nootrealization}. In particular the
Hodge signatures and the commutant are as computed above. Since
\(\dim_{\Q}W=24\), its abelian dimension is twelve. Theorem~\ref{thm:adjoint}
gives the ordinary density assertion.
\end{proof}

In the split CM inequalities of Theorem~\ref{thm:bounds}, these
signatures still permit one \((2,2)\) pair to have slopes
\((0,1/2,1/2,1)\), with every other pair ordinary. For the two mixed
pairs together the rounded sums at \(j=1,2,3\) are \(0,2,4\), all
strictly below \(2j\). The definite pair contributes \(j\) to each
sum. The adjoint criterion gives ordinarity in this example by
using that the mixed coefficient blocks are tensor representations
supported on different embeddings of one rational adjoint factor.
